\chapter*{Список сокращений и обозначений}
\addcontentsline{toc}{chapter}{Список сокращений и обозначений}

\begin{description}
  \item[BP] обратное распространение ошибки.
  \item[PC] предиктивное кодирование.
  \item[Exact] исследуемый режим PC с точным вычислительным отношением,
    определённый реализацией Torch2PC.
  \item[FixedPred] режим PC с фиксированными предсказаниями в пределах
    итеративного вывода состояния.
  \item[Strict] исследуемый строгий режим PC, определённый протоколом проекта.
  \item[CKA] центрированное выравнивание ядер, мера сходства представлений.
  \item[RSA] анализ сходства представлений; в работе используется
    корреляционная форма RSA, зафиксированная протоколом этап~3A.
  \item[VJP] векторно-якобианское произведение.
  \item[ROCm] программная платформа AMD для вычислений на GPU, использованная
    в зафиксированной среде этап~3B.
  \item[CI] непрерывная интеграция; в репозитории также является каноническим
    способом сборки и проверки артефакта диссертации.
  \item[B0] базовая поверхность профилирования этап~3B.
  \item[B1/B2] точные кандидаты реализации этап~3B; прохождение отдельной
    проверки эквивалентности не меняет их роль кандидата до итогового
    инженерного решения.
  \item[SI-MA0/SI-MA1] последовательные этапы отнесения эффекта к механизму и
    калибровки стоимости наблюдателя для итеративного вывода состояния.
  \item[NCZ] зона нулевого вклада: точная тривиальная часть множества
    коррекционного нуля, в которой все канонические каналы равны нулю; в
    численном протоколе отдельно используется зарегистрированная окрестность
    малой активности каналов.
  \item[ECZ] зона компенсации ошибок: точная нетривиальная часть множества
    коррекционного нуля, в которой ненулевые каналы суммируются в нуль; в
    численном протоколе отдельно используется зарегистрированная активная
    окрестность компенсации.
  \item[TNZ] зона нулевого переноса: ненулевой источник ошибки в ядре
    сопряжённого оператора переноса к состоянию, дающий нулевой верхний канал.
  \item[D/U/S] семантика теневых решений QWake-FP: \code{DONE},
    \code{UNKNOWN} и \code{SWEEP}.
  \item[PC-TREF] рамка относительной к задаче эквивалентности
    предиктивного кодирования; связывает диагностическое представление
    состояния с допустимым вычислительным действием.
  \item[PC-CATM] операторная модель агрегации коррекций и переноса ошибок в
    предиктивном кодировании; описывает механизмные источники диагностических
    признаков.
  \item[QWake-PC] собственное имя архитектуры управления вычислительным
    бюджетом, использующей критерий достаточности PC-TREF и механизмные признаки
    PC-CATM; не обозначает один фиксированный алгоритм.
  \item[QWake-FP] ограниченная теневая реализация QWake-PC для
    зарегистрированного случая FixedPred, проверяемая в диссертации.
  \item[C1/C2/C3] последовательные стадии QWake-FP: сбор траекторий и меток
    после действия, офлайн-проверка замороженных правил принятия решения и
    условно открываемый подтверждающий этап.
\end{description}
